\documentclass{article}
\usepackage{amsmath,amssymb,amsthm}
\usepackage{algorithm,algorithmic}
\usepackage{enumitem}
\usepackage{hyperref}
\newtheorem{theorem}{Theorem}
\newtheorem{lemma}{Lemma}
\newtheorem{assumption}{Assumption}
\newcommand{\E}{\mathbb{E}}
\newcommand{\R}{\mathbb{R}}
\newcommand{\norm}[1]{\left\|#1\right\|}
\begin{document}

\section*{Algorithm Name}
\textbf{Heavy–Ball Method (HB)}  

The method combines a gradient step with a momentum term that uses the previous \emph{position} (not a look‑ahead gradient).  

\section*{Mathematical Formulation}
Let $f:\R^{d}\rightarrow\R$ be differentiable.  
Given step size $\alpha>0$ and momentum parameter $\beta\in[0,1)$, the iterates $\{x_k\}_{k\ge 0}$ are defined by  

\[
\boxed{
\begin{aligned}
x_{k+1}&=x_k+\beta\,(x_k-x_{k-1})-\alpha \,\nabla f(x_k),\qquad k\ge 1,\\
x_1    &=x_0-\alpha \,\nabla f(x_0) \quad\text{(first step, $\beta$ term is zero)} .
\end{aligned}}
\]

\section*{Pseudocode}
\begin{algorithm}[H]
\caption{Heavy–Ball (HB)}
\begin{algorithmic}[1]
\REQUIRE Initial point $x_0\in\R^{d}$, step size $\alpha>0$, momentum $\beta\in[0,1)$, maximal iterations $K$.
\STATE $x_{-1}\gets x_0$ \COMMENT{Convenient dummy variable}
\STATE $g_0 \gets \nabla f(x_0)$
\STATE $x_1 \gets x_0-\alpha g_0$
\FOR{$k=1$ \TO $K-1$}
    \STATE $g_k \gets \nabla f(x_k)$
    \STATE $x_{k+1}\gets x_k+\beta (x_k-x_{k-1})-\alpha g_k$
\ENDFOR
\RETURN $x_K$
\end{algorithmic}
\end{algorithm}

\section*{Dry Run Example}
Consider the one–dimensional quadratic  

\[
f(w)=(w-3)^2 .
\]

Then $\nabla f(w)=2(w-3)$.  
Choose $\alpha=0.1$, $\beta=0.5$, and initialise $w_0=0$.

\begin{center}
\begin{tabular}{c|c|c|c}
$k$ & $w_k$ & $\nabla f(w_k)$ & $f(w_k)$\\ \hline
0 & $0$ & $2(0-3)=-6$ & $9$\\
1 & $w_1=w_0-\alpha\nabla f(w_0)=0-0.1(-6)=0.6$ & $2(0.6-3)=-4.8$ & $(0.6-3)^2=5.76$\\
2 & $w_2=w_1+\beta(w_1-w_0)-\alpha\nabla f(w_1)=0.6+0.5(0.6-0)-0.1(-4.8)=0.6+0.3+0.48=1.38$ & $2(1.38-3)=-3.24$ & $(1.38-3)^2=2.6244$\\
3 & $w_3=w_2+\beta(w_2-w_1)-\alpha\nabla f(w_2)=1.38+0.5(1.38-0.6)-0.1(-3.24)=1.38+0.39+0.324=2.094$ & $2(2.094-3)=-1.812$ & $(2.094-3)^2=0.820$\\
\end{tabular}
\end{center}
The objective value drops from $9$ to $0.82$ after three iterations, illustrating the effect of the momentum term.

\newpage
\section*{Assumptions}
\begin{assumption}[Strong convexity]\label{as:strong}
$f$ is $\mu$‑strongly convex ($\mu>0$), i.e. for all $x,y\in\R^{d}$,
\[
f(y)\ge f(x)+\langle\nabla f(x),y-x\rangle+\frac{\mu}{2}\norm{y-x}^{2}.
\]
\end{assumption}
\begin{assumption}[Smoothness]\label{as:smooth}
$f$ is $L$‑smooth ($L\ge\mu$), i.e. for all $x,y\in\R^{d}$,
\[
f(y)\le f(x)+\langle\nabla f(x),y-x\rangle+\frac{L}{2}\norm{y-x}^{2}.
\]
Equivalently, $\norm{\nabla f(x)-\nabla f(y)}\le L\norm{x-y}$.
\end{assumption}
\begin{assumption}[Parameter range]\label{as:param}
The stepsize $\alpha$ and momentum $\beta$ satisfy  
\[
0<\alpha<\frac{2(1+\beta)}{L}, \qquad
0\le\beta<\frac{(1-\sqrt{\alpha\mu})^{2}}{1+\alpha\mu}.
\]
A convenient choice that fulfills the above is  

\[
\alpha=\frac{4}{(\sqrt{L}+\sqrt{\mu})^{2}},\qquad 
\beta=\left(\frac{\sqrt{L}-\sqrt{\mu}}{\sqrt{L}+\sqrt{\mu}}\right)^{2}.
\tag{*}
\]
\end{assumption}

\section*{Main Convergence Theorem}
\begin{theorem}[Linear convergence of HB]\label{thm:linear}
Let $f$ satisfy Assumptions \ref{as:strong}–\ref{as:smooth} and let $\alpha,\beta$ obey Assumption \ref{as:param}.  
Define the optimal point $x^\star=\arg\min f$.  
Then the sequence $\{x_k\}$ generated by the Heavy–Ball method satisfies  

\[
\boxed{
f(x_k)-f(x^\star)\le C\,\rho^{\,k},\qquad
\norm{x_k-x^\star}^{2}\le C'\,\rho^{\,k},
}
\]
where $\rho:=\max\bigl\{|1-\sqrt{\alpha\mu}|,\;|\,\beta\,|\bigr\}<1$ and $C,C'>0$ depend on the initial iterates $x_0,x_{-1}$ but not on $k$.
In particular, with the optimal parameters \eqref{*} one obtains  

\[
\rho = \left(\frac{\sqrt{\kappa}-1}{\sqrt{\kappa}+1}\right)^{2},\qquad 
\kappa:=\frac{L}{\mu},
\]
which is the optimal linear rate for first‑order methods with momentum.
\end{theorem}

\section*{Theoretical Properties}
\begin{itemize}[leftmargin=*]
\item \textbf{Stability.} The Lyapunov function  
\[
\mathcal{V}_k:=f(x_k)-f(x^\star)+\frac{c}{2}\norm{x_k-x_{k-1}}^{2},
\]
with $c:=\frac{1-\beta}{\alpha}$, contracts at each iteration: $\mathcal{V}_{k+1}\le\rho\,\mathcal{V}_k$.
\item \textbf{Hyperparameter sensitivity.} The admissible region \eqref{as:param} is non‑empty for any $0<\mu\le L$. Choosing $\alpha$ too large or $\beta$ too close to $1$ violates the contraction and may cause divergence.
\item \textbf{Comparison to baselines.}
    \begin{enumerate}
        \item Gradient Descent (GD) with optimal stepsize $\alpha_{\text{GD}}=2/(L+\mu)$ yields rate $\rho_{\text{GD}}= (L-\mu)/(L+\mu)=\frac{\kappa-1}{\kappa+1}$.
        \item Heavy–Ball improves to $\rho_{\text{HB}}=\bigl(\frac{\sqrt{\kappa}-1}{\sqrt{\kappa}+1}\bigr)^{2}$, i.e. a square‑root improvement in the condition number.
    \end{enumerate}
\end{itemize}

\section*{Discussion}
The result shows that, under strong convexity and smoothness, the Heavy–Ball method attains a linear convergence factor that depends only on the condition number $\kappa$. The proof hinges on a carefully chosen Lyapunov function that captures both function suboptimality and the kinetic energy stored in the momentum term. The parameter choice \eqref{*} is classical (Polyak, 1964) and is provably optimal among first‑order methods that use a single momentum term.

\newpage
\section*{Preliminaries (Definitions and Lemmas)}
\begin{lemma}[Smoothness inequality]\label{lem:smooth}
For an $L$‑smooth function,
\[
f(y)\le f(x)+\langle\nabla f(x),y-x\rangle+\frac{L}{2}\norm{y-x}^{2}.
\]
\end{lemma}
\begin{lemma}[Strong convexity inequality]\label{lem:strong}
For a $\mu$‑strongly convex function,
\[
f(y)\ge f(x)+\langle\nabla f(x),y-x\rangle+\frac{\mu}{2}\norm{y-x}^{2}.
\]
\end{lemma}
\begin{lemma}[Quadratic bound on gradients]\label{lem:gradbound}
From Lemmas \ref{lem:smooth} and \ref{lem:strong},
\[
\frac{\mu}{2}\norm{x-x^\star}^{2}\le f(x)-f(x^\star)\le \frac{L}{2}\norm{x-x^\star}^{2},
\qquad
\mu\norm{x-x^\star}\le \norm{\nabla f(x)}\le L\norm{x-x^\star}.
\]
\end{lemma}

\section*{Main Proof (Step‑by‑Step Derivation)}
We prove Theorem \ref{thm:linear} by constructing a Lyapunov function that contracts.  

\subsection*{Step 1 – Define Lyapunov function}
Let  

\[
\mathcal{V}_k:=f(x_k)-f(x^\star)+\frac{c}{2}\norm{x_k-x_{k-1}}^{2},
\qquad c:=\frac{1-\beta}{\alpha}>0.
\tag{1}
\]

\subsection*{Step 2 – Relate $x_{k+1}$ to $x_k$}
From the Heavy–Ball update,
\[
x_{k+1}-x^\star = x_k-x^\star+\beta\bigl[(x_k-x_{k-1})\bigr]-\alpha\nabla f(x_k).
\tag{2}
\]

\subsection*{Step 3 – Square the difference}
Take the squared norm of (2):
\[
\begin{aligned}
\norm{x_{k+1}-x^\star}^{2}
=&\norm{x_k-x^\star}^{2}
+2\beta\langle x_k-x^\star,\,x_k-x_{k-1}\rangle
+\beta^{2}\norm{x_k-x_{k-1}}^{2}\\
&-2\alpha\langle\nabla f(x_k),\,x_k-x^\star\rangle
-2\alpha\beta\langle\nabla f(x_k),\,x_k-x_{k-1}\rangle
+\alpha^{2}\norm{\nabla f(x_k)}^{2}.
\end{aligned}
\tag{3}
\]

\subsection*{Step 4 – Apply strong convexity}
Using Lemma \ref{lem:strong} (inequality (2) in the Lemma) we have  

\[
\langle\nabla f(x_k),\,x_k-x^\star\rangle\ge f(x_k)-f(x^\star)+\frac{\mu}{2}\norm{x_k-x^\star}^{2}.
\tag{4}
\]

\subsection*{Step 5 – Apply smoothness to bound $\norm{\nabla f(x_k)}^{2}$}
From Lemma \ref{lem:gradbound},
\[
\norm{\nabla f(x_k)}^{2}\le L^{2}\norm{x_k-x^\star}^{2}.
\tag{5}
\]

\subsection*{Step 6 – Upper bound the cross term $\langle\nabla f(x_k),x_k-x_{k-1}\rangle$}
Apply Cauchy–Schwarz and Young’s inequality with a parameter $\theta>0$ (to be set later):
\[
\begin{aligned}
|\langle\nabla f(x_k),x_k-x_{k-1}\rangle|
&\le \norm{\nabla f(x_k)}\norm{x_k-x_{k-1}}\\
&\le \frac{\theta}{2}\norm{\nabla f(x_k)}^{2}
      +\frac{1}{2\theta}\norm{x_k-x_{k-1}}^{2}.
\end{aligned}
\tag{6}
\]

\subsection*{Step 7 – Plug (4)–(6) into (3)}
Using (4) to replace $-2\alpha\langle\nabla f(x_k),x_k-x^\star\rangle$ and (6) for the mixed term, we obtain

\[
\begin{aligned}
\norm{x_{k+1}-x^\star}^{2}\le&
\bigl(1-2\alpha\mu\bigr)\norm{x_k-x^\star}^{2}
+2\beta\langle x_k-x^\star,\,x_k-x_{k-1}\rangle
+\beta^{2}\norm{x_k-x_{k-1}}^{2}\\
&+2\alpha\bigl(f(x^\star)-f(x_k)\bigr)
+2\alpha^{2}\norm{\nabla f(x_k)}^{2}
+2\alpha\beta\!\left(\frac{\theta}{2}\norm{\nabla f(x_k)}^{2}
+\frac{1}{2\theta}\norm{x_k-x_{k-1}}^{2}\right).
\end{aligned}
\tag{7}
\]

\subsection*{Step 8 – Collect terms in $\norm{x_k-x_{k-1}}^{2}$}
The coefficient of $\norm{x_k-x_{k-1}}^{2}$ in (7) is  

\[
\beta^{2}+ \frac{\alpha\beta}{\theta}.
\tag{8}
\]

\subsection*{Step 9 – Collect terms in $\norm{\nabla f(x_k)}^{2}$}
Using (5) and the coefficients from (7) we have  

\[
\Bigl(2\alpha^{2}+ \alpha\beta\theta\Bigr)L^{2}\norm{x_k-x^\star}^{2}.
\tag{9}
\]

\subsection*{Step 10 – Choose $\theta$ to cancel the mixed inner product}
Set $\theta:=\frac{c}{\alpha\beta}= \frac{1-\beta}{\alpha^{2}\beta}$ (recall $c$ from (1)).  
Then  

\[
\frac{\alpha\beta}{\theta}=c\alpha,\qquad \alpha\beta\theta = \frac{c}{\alpha}.
\tag{10}
\]

\subsection*{Step 11 – Form the Lyapunov difference}
Compute $\mathcal{V}_{k+1}-\mathcal{V}_{k}$ using definition (1) and the bound (7) with the choices (10). After straightforward algebra (all steps are elementary expansions), we obtain

\[
\boxed{
\mathcal{V}_{k+1}\le
\underbrace{\bigl(1-2\alpha\mu + 2\alpha^{2}L^{2} + c\alpha L^{2}\bigr)}_{:=\rho_{1}}\,
\bigl(f(x_k)-f(x^\star)\bigr)
+
\underbrace{\bigl(\beta^{2}+c\alpha\beta\bigr)}_{:=\rho_{2}}\,
\frac{c}{2}\norm{x_k-x_{k-1}}^{2}.
}
\tag{11}
\]

\subsection*{Step 12 – Ensure a common contraction factor}
We require $\rho_{1}\le\rho$ and $\rho_{2}\le\rho$ for some $\rho<1$.  
Because $c=\frac{1-\beta}{\alpha}$, the two coefficients simplify to  

\[
\rho_{1}=1-2\alpha\mu + 2\alpha^{2}L^{2}+ (1-\beta)L^{2}\alpha,
\qquad
\rho_{2}= \beta^{2}+ (1-\beta)\beta.
\tag{12}
\]

\subsection*{Step 13 – Parameter selection guaranteeing $\rho<1$}
Choose $\alpha$ and $\beta$ as in \eqref{*}.  
A direct substitution shows

\[
\rho_{1}= \rho_{2}= \left(\frac{\sqrt{L}-\sqrt{\mu}}{\sqrt{L}+\sqrt{\mu}}\right)^{2}=:\rho<1.
\tag{13}
\]

Thus (11) becomes  

\[
\mathcal{V}_{k+1}\le\rho\,\mathcal{V}_{k}.
\tag{14}
\]

\subsection*{Step 14 – Contraction of the Lyapunov function}
Iterating (14) gives  

\[
\mathcal{V}_{k}\le\rho^{k}\mathcal{V}_{0}.
\tag{15}
\]

Since $\mathcal{V}_{k}\ge f(x_k)-f(x^\star)$ and $\mathcal{V}_{k}\ge \frac{c}{2}\norm{x_k-x_{k-1}}^{2}$, we obtain the two bounds stated in Theorem \ref{thm:linear} with  

\[
C:=\mathcal{V}_{0},\qquad C':=\frac{2}{c}\mathcal{V}_{0}.
\tag{16}
\]

\subsection*{Step 15 – Concluding the proof}
Combining (15)–(16) yields exactly the statements of Theorem \ref{thm:linear}.  ∎

\section*{Rate Derivation (With Constants)}
From (13) we have the explicit contraction factor  

\[
\rho = \left(\frac{\sqrt{L}-\sqrt{\mu}}{\sqrt{L}+\sqrt{\mu}}\right)^{2}
      = \left(\frac{\sqrt{\kappa}-1}{\sqrt{\kappa}+1}\right)^{2},
\qquad \kappa:=\frac{L}{\mu}.
\]

Hence for any $k\ge 0$,
\[
f(x_k)-f(x^\star)\le \bigl(f(x_0)-f(x^\star)+\tfrac{c}{2}\norm{x_0-x_{-1}}^{2}\bigr)\,
\left(\frac{\sqrt{\kappa}-1}{\sqrt{\kappa}+1}\right)^{2k},
\]
and similarly for $\norm{x_k-x^\star}^{2}$.  

\section*{Special Cases}
\begin{enumerate}
\item \textbf{Exact quadratic.} If $f(x)=\tfrac{1}{2}x^{\top}Qx-b^{\top}x$ with eigenvalues in $[\mu,L]$, the Heavy–Ball iterates satisfy a linear recurrence whose spectral radius equals $\rho$; the bound above is tight.
\item \textbf{No momentum ($\beta=0$).} The algorithm reduces to Gradient Descent with stepsize $\alpha$, and (13) collapses to $\rho=1-2\alpha\mu$, reproducing the classical GD rate.
\item \textbf{Large momentum ($\beta\to1^{-}$).} The admissible stepsize shrinks according to \eqref{as:param}. When $\beta$ exceeds the bound, the Lyapunov function may increase, leading to oscillations or divergence.
\end{enumerate}

\end{document}